\documentclass[aps,preprint]{revtex4}%
\usepackage{amsfonts}%
\usepackage{amsmath}%
\setcounter{MaxMatrixCols}{30}%
\usepackage{amssymb}%
\usepackage{graphicx}
%TCIDATA{OutputFilter=latex2.dll}
%TCIDATA{Version=4.10.0.2363}
%TCIDATA{CSTFile=revtex4.cst}
%TCIDATA{Created=Friday, November 28, 2003 11:55:12}
%TCIDATA{LastRevised=Friday, November 28, 2003 11:55:37}
%TCIDATA{<META NAME="GraphicsSave" CONTENT="32">}
%TCIDATA{<META NAME="DocumentShell" CONTENT="Articles\SW\REVTeX 4">}
\newtheorem{theorem}{Theorem}
\newtheorem{acknowledgement}[theorem]{Acknowledgement}
\newtheorem{algorithm}[theorem]{Algorithm}
\newtheorem{axiom}[theorem]{Axiom}
\newtheorem{claim}[theorem]{Claim}
\newtheorem{conclusion}[theorem]{Conclusion}
\newtheorem{condition}[theorem]{Condition}
\newtheorem{conjecture}[theorem]{Conjecture}
\newtheorem{corollary}[theorem]{Corollary}
\newtheorem{criterion}[theorem]{Criterion}
\newtheorem{definition}[theorem]{Definition}
\newtheorem{example}[theorem]{Example}
\newtheorem{exercise}[theorem]{Exercise}
\newtheorem{lemma}[theorem]{Lemma}
\newtheorem{notation}[theorem]{Notation}
\newtheorem{problem}[theorem]{Problem}
\newtheorem{proposition}[theorem]{Proposition}
\newtheorem{remark}[theorem]{Remark}
\newtheorem{solution}[theorem]{Solution}
\newtheorem{summary}[theorem]{Summary}
\newenvironment{proof}[1][Proof]{\noindent\textbf{#1.} }{\ \rule{0.5em}{0.5em}}
\begin{document}
\preprint{HEP/123-qed}
\title[Short title for running header]{Long title that appears on the title page}
\author{First Author}
\affiliation{My Institution}
\author{Second Author}
\affiliation{My Institution}
\author{Third Author}
\affiliation{Other Institution}
\keywords{one two three}
\pacs{PACS number}

\begin{abstract}
Shell document for REV\TeX{} 4.

\end{abstract}
\volumeyear{year}
\volumenumber{number}
\issuenumber{number}
\eid{identifier}
\date[Date text]{date}
\received[Received text]{date}

\revised[Revised text]{date}

\accepted[Accepted text]{date}

\published[Published text]{date}

\startpage{101}
\endpage{102}
\maketitle
\tableofcontents


\section{}

The perturbing potential $\mathcal{F}\left(  \mathrm{z,\bar{z}}\mathbf{,}%
t\right)  $ can be expanded in terms of the incremental paths $\left\{
\delta\mathrm{z}\left(  t\right)  \right\}  $ about a critical (stationary)
path $\mathrm{z}_{0}\left(  t\right)  =\mathbf{z}_{0}$ as
\begin{align}
\mathcal{F}\left(  \mathbf{z}_{0}+\delta\mathrm{z}\left(  t\right)  ,\bar
{z}_{0}\left(  t\right)  +\delta\mathrm{\bar{z}}\left(  t\right)
\mathbf{,}t\right)   &  =\mathcal{F}\left(  z_{0},\bar{z}_{0}\mathbf{,}%
t\right)  +D_{\mathbf{z}}\mathcal{F}\left(  z_{0},\bar{z}_{0}\mathbf{,}%
t\right)  \bullet\delta\mathrm{z}\left(  t\right)  \nonumber\\
&  +D_{\mathbf{\bar{z}}}\mathcal{F}\left(  z_{0},\bar{z}_{0}\mathbf{,}%
t\right)  \bullet\delta\mathrm{\bar{z}}\left(  t\right)  +\cdots
\end{align}
where
\begin{equation}
D_{\mathbf{x}}\mathcal{F}\left(  \mathbf{z}_{0},\mathbf{\bar{z}}_{0}%
\mathbf{,}t\right)  =D_{\mathbf{\rho}}\mathcal{F}\left(  \mathbf{z}%
_{0},\mathbf{\bar{z}}_{0}\mathbf{,}t\right)  \bullet D_{\mathbf{x}%
}\mathcal{\rho}\left(  \mathbf{z}_{0},\mathbf{\bar{z}}_{0}\right)  =\sum
_{k,l}f_{kl}\left(  t\right)  \frac{\partial\mathcal{\rho}_{kl}\left(
\mathbf{z}_{0},\mathbf{\bar{z}}_{0}\right)  }{\partial x_{j}};\quad
\mathbf{x=z,\bar{z}}%
\end{equation}
and the matrix elements of $\left\{  D_{\mathbf{z}}\mathcal{F}\left(
\mathbf{z}_{0},\mathbf{\bar{z}}_{0}\mathbf{,}t\right)  ,D_{\mathbf{\rho}%
}\mathcal{F}\left(  \mathbf{z}_{0},\mathbf{\bar{z}}_{0}\mathbf{,}t\right)
,D_{\mathbf{z}}\mathcal{\rho}\left(  \mathbf{z}_{0},\mathbf{\bar{z}}%
_{0}\right)  \right\}  $ are given by
\begin{align}
\left[  D_{\mathbf{x}}^{1}\mathcal{F}\left(  \mathbf{z}_{0},\mathbf{\bar{z}%
}_{0}\mathbf{,}t\right)  \right]  _{j} &  =\frac{\partial\mathcal{F}\left(
\mathbf{z}_{0},\mathbf{\bar{z}}_{0}\mathbf{,}t\right)  }{\partial x_{j}}%
=\sum_{k,l}\frac{\partial\mathcal{F}\left(  \mathbf{z}_{0},\mathbf{\bar{z}%
}_{0}\mathbf{,}t\right)  }{\partial\rho_{kl}}\frac{\partial\mathcal{\rho}%
_{kl}\left(  \mathbf{z}_{0},\mathbf{\bar{z}}_{0}\right)  }{\partial z_{j}%
};\mathbf{x=z,\bar{z};}x_{j}=z_{j},\bar{z}_{j}\nonumber\\
\left[  D_{\mathbf{\rho}}^{1}\mathcal{F}\left(  \mathbf{z}_{0},\mathbf{\bar
{z}}_{0}\mathbf{,}t\right)  \right]  _{kl} &  =\frac{\partial\mathcal{F}%
\left(  \mathbf{z}_{0},\mathbf{\bar{z}}_{0}\mathbf{,}t\right)  }{\partial
\rho_{kl}}=f_{kl}\left(  t\right)  \nonumber\\
\left[  D_{\mathbf{x}}^{1}\mathcal{\rho}\left(  \mathbf{z}_{0},\mathbf{\bar
{z}}_{0}\right)  \right]  _{klj} &  =\frac{\partial\mathcal{\rho}_{kl}\left(
\mathbf{z}_{0},\mathbf{\bar{z}}_{0}\right)  }{\partial x_{j}};\mathbf{x=z,\bar
{z};}x_{j}=z_{j},\bar{z}_{j}%
\end{align}
leading to
\begin{align}
\delta_{\mathbf{z}_{0}}^{\left(  1\right)  }\mathcal{F}\left(  \delta
\mathrm{z}\left(  t\right)  ,\delta\mathrm{\bar{z}}\left(  t\right)
\mathbf{,}t\right)   &  =\sum_{k,l}f_{kl}\left(  t\right)  \delta^{\left(
1\right)  }\rho_{lk}\left(  \delta\mathrm{z}\left(  t\right)  ,\delta
\mathrm{\bar{z}}\left(  t\right)  \right)  \nonumber\\
&  =\sum_{k,l}f_{kl}\left(  t\right)  \left\{  \frac{\partial\mathcal{\rho
}_{kl}\left(  \mathbf{z}_{0},\mathbf{\bar{z}}_{0}\right)  }{\partial z_{j}%
}\delta\mathrm{z}_{j}\left(  t\right)  +\frac{\partial\mathcal{\rho}%
_{kl}\left(  \mathbf{z}_{0},\mathbf{\bar{z}}_{0}\right)  }{\partial\bar{z}%
_{j}}\delta\mathrm{\bar{z}}_{j}\left(  t\right)  \right\}
\end{align}
One can form a quadratic approximation to this evolution around a point
$\left(  \mathbf{z}_{0},\mathbf{\bar{z}}_{0}\right)  $ by expanding up to
second order in the local changes $\left(  \delta\mathbf{z}\left(  t\right)
,\delta\mathbf{\bar{z}}\left(  t\right)  \right)  $
\begin{equation}
\mathcal{E}\left(  \mathbf{z},\mathbf{\bar{z},}t\right)  \approx
\mathcal{E}_{2}\left[  \mathbf{z}_{0}\right]  \left(  \delta\mathbf{z}%
,\delta\mathbf{\bar{z}},t\right)  =\mathcal{E}\left(  \mathbf{z}%
_{0},\mathbf{\bar{z}}_{0}\mathbf{,}t\right)  +\delta^{1}\mathcal{E}\left(
\mathbf{z}_{0},\mathbf{\bar{z}}_{0}\mathbf{,}t\right)  +\delta^{2}%
\mathcal{E}\left(  \mathbf{z}_{0},\mathbf{\bar{z}}_{0}\mathbf{,}t\right)
\end{equation}
to produce an approximate classical Hamiltonian around the point $\left(
\mathbf{z}_{0},\mathbf{\bar{z}}_{0}\right)  $ where
\begin{align}
\delta^{1}\mathcal{E}\left(  \mathbf{z}_{0},\mathbf{\bar{z}}_{0}%
\mathbf{,}t\right)   &  =D^{1}\mathcal{E}\left(  \mathbf{z}_{0},\mathbf{\bar
{z}}_{0},t\right)  \left(
\begin{array}
[c]{c}%
\delta\mathbf{z}\left(  t\right)  \\
\delta\mathbf{\bar{z}}\left(  t\right)
\end{array}
\right)  \\
\delta^{2}\mathcal{E}\left(  \mathbf{z}_{0},\mathbf{\bar{z}}_{0}%
\mathbf{,}t\right)   &  =\frac{1}{2}\left(
\begin{array}
[c]{cc}%
\delta\mathbf{z}\left(  t\right)   & \delta\mathbf{\bar{z}}\left(  t\right)
\end{array}
\right)  D^{2}\mathcal{E}\left(  \mathbf{z}_{0},\mathbf{\bar{z}}_{0},t\right)
\left(
\begin{array}
[c]{c}%
\delta\mathbf{z}\left(  t\right)  \\
\delta\mathbf{\bar{z}}\left(  t\right)
\end{array}
\right)
\end{align}
The gradient $D^{1}\mathcal{E}\left(  \mathbf{z}_{0},\mathbf{\bar{z}}%
_{0},t\right)  $ of $\mathcal{E}\left(  \mathbf{z},\mathbf{\bar{z}},t\right)
$ at the critical point $\left(  \mathbf{z}_{0},\mathbf{\bar{z}}_{0}\right)  $
of $\mathcal{E}_{0}$ is given by
\begin{align*}
\frac{\partial\mathcal{E}\left(  \mathbf{z}_{0},\mathbf{\bar{z}}_{0},t\right)
}{\partial\mathbf{x}} &  =\frac{1}{\mathcal{S}\left(  \mathbf{z}%
_{0},\mathbf{\bar{z}}_{0}\right)  }\left\{  \frac{\partial\mathcal{H}%
_{0}\left(  \mathbf{z}_{0},\mathbf{\bar{z}}_{0}\right)  }{\partial\mathbf{x}%
}+\frac{\partial\mathcal{F}\left(  \mathbf{z}_{0},\mathbf{\bar{z}}%
_{0},t\right)  }{\partial\mathbf{x}}\right.  \\
&  \qquad\qquad\qquad\qquad\qquad\left.  -\left(  \mathcal{E}_{0}\left(
\mathbf{z}_{0},\mathbf{\bar{z}}_{0}\right)  +\mathcal{F}\left(  \mathbf{z}%
_{0},\mathbf{\bar{z}}_{0},t\right)  \right)  \frac{\partial\mathcal{S}\left(
\mathbf{z}_{0},\mathbf{\bar{z}}_{0}\right)  }{\partial\mathbf{x}}\right\}  \\
&  =\frac{1}{\mathcal{S}\left(  \mathbf{z}_{0},\mathbf{\bar{z}}_{0}\right)
}\left\{  \frac{\partial\mathcal{F}\left(  \mathbf{z}_{0},\mathbf{\bar{z}}%
_{0},t\right)  }{\partial\mathbf{x}}-\mathcal{F}\left(  \mathbf{z}%
_{0},\mathbf{\bar{z}}_{0},t\right)  \frac{\partial\mathcal{S}\left(
\mathbf{z}_{0},\mathbf{\bar{z}}_{0}\right)  }{\partial\mathbf{x}}\right\}
\end{align*}
as
\begin{equation}
\frac{\partial\mathcal{H}_{0}\left(  \mathbf{z}_{0},\mathbf{\bar{z}}%
_{0}\right)  }{\partial\mathbf{x}}-\mathcal{E}_{0}\left(  \mathbf{z}%
_{0},\mathbf{\bar{z}}_{0}\right)  \frac{\partial\mathcal{S}\left(
\mathbf{z}_{0},\mathbf{\bar{z}}_{0}\right)  }{\partial\mathbf{x}}=\mathbf{0}%
\end{equation}
The Hessian $D^{2}\mathcal{E}\left(  \mathbf{z}_{0},\mathbf{\bar{z}}%
_{0},t\right)  $ at any point $\left(  \mathbf{z}_{0},\mathbf{\bar{z}}%
_{0},t\right)  $ is given by matrix elements
\begin{align}
&  \frac{\partial^{2}\mathcal{E}\left(  \mathbf{z}_{0},\mathbf{\bar{z}}%
_{0},t\right)  }{\partial\mathbf{x}\partial\mathbf{y}}=\frac{1}{\mathcal{S}%
\left(  \mathbf{z}_{0},\mathbf{\bar{z}}_{0},t\right)  }\left\{  \frac
{\partial^{2}\mathcal{H}\left(  \mathbf{z}_{0},\mathbf{\bar{z}}_{0},t\right)
}{\partial\mathbf{x}\partial\mathbf{y}}-\frac{\partial\mathcal{E}\left(
\mathbf{z}_{0},\mathbf{\bar{z}}_{0},t\right)  }{\partial\mathbf{x}}%
\frac{\partial\mathcal{S}\left(  \mathbf{z}_{0},\mathbf{\bar{z}}_{0},t\right)
}{\partial\mathbf{y}}\right.  \nonumber\\
&  \qquad\qquad\qquad\qquad\qquad\qquad\left.  -\mathcal{E}\left(
\mathbf{z}_{0},\mathbf{\bar{z}}_{0},t\right)  \frac{\partial^{2}%
\mathcal{S}\left(  \mathbf{z}_{0},\mathbf{\bar{z}}_{0},t\right)  }%
{\partial\mathbf{x}\partial\mathbf{y}}-\frac{\partial\mathcal{S}\left(
\mathbf{z}_{0},\mathbf{\bar{z}}_{0},t\right)  }{\partial\mathbf{x}}%
\frac{\partial\mathcal{E}\left(  \mathbf{z}_{0},\mathbf{\bar{z}}_{0},t\right)
}{\partial\mathbf{y}}\right\}  \nonumber\\
&  \qquad\qquad\qquad\quad\left.  =\right.  \frac{1}{\mathcal{S}\left(
\mathbf{z}_{0},\mathbf{\bar{z}}_{0},t\right)  }\left\{  \frac{\partial
^{2}\mathcal{H}_{0}\left(  \mathbf{z}_{0},\mathbf{\bar{z}}_{0},t\right)
}{\partial\mathbf{x}\partial\mathbf{y}}+\frac{\partial^{2}\mathcal{F}\left(
\mathbf{z}_{0},\mathbf{\bar{z}}_{0},t\right)  }{\partial\mathbf{x}%
\partial\mathbf{y}}\right.  \nonumber\\
&  -\left(  \frac{\partial\mathcal{E}_{0}\left(  \mathbf{z}_{0},\mathbf{\bar
{z}}_{0},t\right)  }{\partial\mathbf{x}}+\frac{\partial\mathcal{F}\left(
\mathbf{z}_{0},\mathbf{\bar{z}}_{0},t\right)  }{\partial\mathbf{x}}\right)
\frac{\partial\mathcal{S}\left(  \mathbf{z}_{0},\mathbf{\bar{z}}_{0},t\right)
}{\partial\mathbf{y}}-\left(  \mathcal{E}_{0}\left(  \mathbf{z}_{0}%
,\mathbf{\bar{z}}_{0}\right)  +\mathcal{F}\left(  \mathbf{z}_{0}%
,\mathbf{\bar{z}}_{0},t\right)  \right)  \frac{\partial^{2}\mathcal{S}\left(
\mathbf{z}_{0},\mathbf{\bar{z}}_{0}\right)  }{\partial\mathbf{x}%
\partial\mathbf{y}}\nonumber\\
&  \qquad\qquad\qquad\qquad\qquad\qquad\qquad\qquad\qquad\qquad\left.
-\frac{\partial\mathcal{S}\left(  \mathbf{z}_{0},\mathbf{\bar{z}}%
_{0},t\right)  }{\partial\mathbf{x}}\left(  \frac{\partial\mathcal{E}%
_{0}\left(  \mathbf{z}_{0},\mathbf{\bar{z}}_{0}\right)  }{\partial\mathbf{y}%
}+\frac{\partial\mathcal{F}\left(  \mathbf{z}_{0},\mathbf{\bar{z}}%
_{0},t\right)  }{\partial\mathbf{y}}\right)  \right\}
\end{align}
If the point $\left(  \mathbf{z}_{0},\mathbf{\bar{z}}_{0}\right)  $ is a
critical point of the unperturbed energy $\mathcal{E}_{0}\left(
\mathbf{z}_{0},\mathbf{\bar{z}}_{0}\right)  $ and we only consider first order
effects of $\mathcal{F},$ then this gives
\begin{align}
\frac{\partial^{2}\mathcal{E}\left(  \mathbf{z}_{0},\mathbf{\bar{z}}%
_{0},t\right)  }{\partial\mathbf{x}\partial\mathbf{y}} &  =\frac
{1}{\mathcal{S}\left(  \mathbf{z}_{0},\mathbf{\bar{z}}_{0}\right)  }\left\{
\frac{\partial^{2}\mathcal{H}_{0}\left(  \mathbf{z}_{0},\mathbf{\bar{z}}%
_{0}\right)  }{\partial\mathbf{x}\partial\mathbf{y}}\right.  -\frac
{\partial\mathcal{F}\left(  \mathbf{z}_{0},\mathbf{\bar{z}}_{0},t\right)
}{\partial\mathbf{x}}\frac{\partial\mathcal{S}\left(  \mathbf{z}%
_{0},\mathbf{\bar{z}}_{0},t\right)  }{\partial\mathbf{y}}\nonumber\\
&  \left.  -\left(  \mathcal{E}_{0}\left(  \mathbf{z}_{0},\mathbf{\bar{z}}%
_{0}\right)  +\mathcal{F}\left(  \mathbf{z}_{0},\mathbf{\bar{z}}_{0},t\right)
\right)  \frac{\partial^{2}\mathcal{S}\left(  \mathbf{z}_{0},\mathbf{\bar{z}%
}_{0}\right)  }{\partial\mathbf{x}\partial\mathbf{y}}-\frac{\partial
\mathcal{S}\left(  \mathbf{z}_{0},\mathbf{\bar{z}}_{0},t\right)  }%
{\partial\mathbf{x}}\frac{\partial\mathcal{F}\left(  \mathbf{z}_{0}%
,\mathbf{\bar{z}}_{0},t\right)  }{\partial\mathbf{y}}\right\}
\end{align}
A linearized EOM in the tangent space at the point $\left(  \mathbf{z}%
_{0},\mathbf{\bar{z}}_{0}\right)  $ is obtained by using the quadratic
approximation $\mathcal{E}_{2}\left[  \mathbf{z}_{0}\right]  \left(
\delta\mathbf{z},\delta\mathbf{\bar{z}},t\right)  $ to the classical
Hamiltonian around the critical point $\left(  \mathbf{z}_{0},\mathbf{\bar{z}%
}_{0}\right)  $%
\begin{align}
&  \left(
\begin{array}
[c]{c}%
\frac{d\delta\mathbf{\bar{z}}\left(  t\right)  }{dt}\\
\frac{d\delta\mathbf{z}\left(  t\right)  }{dt}%
\end{array}
\right)  \nonumber\\
&  \quad=\left(
\begin{array}
[c]{cc}%
\mathbf{0} & \frac{i}{\hslash}\mathbf{\bar{C}}^{-1}\left(  \mathbf{z}%
_{0},\mathbf{\bar{z}}_{0},t\right)  \\
-\frac{i}{\hslash}\mathbf{C}^{-1}\left(  \mathbf{z}_{0},\mathbf{\bar{z}}%
_{0},t\right)   & \mathbf{0}%
\end{array}
\right)  \left(  D^{1}\mathcal{E}\left(  \mathbf{z}_{0},\mathbf{\bar{z}}%
_{0},t\right)  +D^{2}\mathcal{E}\left(  \mathbf{z}_{0},\mathbf{\bar{z}}%
_{0},t\right)  \left(
\begin{array}
[c]{c}%
\delta\mathbf{\bar{z}}\left(  t\right)  \\
\delta\mathbf{z}\left(  t\right)
\end{array}
\right)  \right)
\end{align}
where the driving force is given by the gradient of $\mathcal{E}_{2}\left[
\mathbf{z}_{0}\right]  $
\begin{equation}
\left(
\begin{array}
[c]{c}%
\frac{\partial\mathcal{E}_{2}\left[  \mathbf{z}_{0}\right]  \left(
\delta\mathbf{z},\delta\mathbf{\bar{z}},t\right)  }{\partial\delta
\mathbf{\bar{z}}}\\
\frac{\partial\mathcal{E}_{2}\left[  \mathbf{z}_{0}\right]  \left(
\delta\mathbf{z},\delta\mathbf{\bar{z}},t\right)  }{\partial\delta\mathbf{z}}%
\end{array}
\right)  =D^{1}\mathcal{E}\left(  \mathbf{z}_{0},\mathbf{\bar{z}}%
_{0},t\right)  +D^{2}\mathcal{E}\left(  \mathbf{z}_{0},\mathbf{\bar{z}}%
_{0},t\right)  \left(
\begin{array}
[c]{c}%
\delta\mathbf{\bar{z}}\left(  t\right)  \\
\delta\mathbf{z}\left(  t\right)
\end{array}
\right)
\end{equation}


The linear response $\delta\mathbf{\rho}\left(  t\right)  $ of $\mathbf{\rho
}\left(  \mathbf{z}_{0},\mathbf{\bar{z}}_{0},t\right)  $ is given by the
linear response $\left\{  \delta\rho_{kl}\left(  t\right)  \right\}  $ of the
classical functions $\rho_{kl}\left(  \mathbf{z}_{0},\mathbf{\bar{z}}%
_{0},t\right)  $
\begin{equation}
\delta\rho_{kl}\left(  t\right)  =\sum_{j}\left\{  \frac{\partial\rho
_{kl}\left(  \mathbf{z}_{0},\mathbf{\bar{z}}_{0},t\right)  }{\partial z_{j}%
}\delta z_{j}\left(  t\right)  +\frac{\partial\rho_{kl}\left(  \mathbf{z}%
_{0},\mathbf{\bar{z}}_{0},t\right)  }{\partial\bar{z}_{j}}\delta\bar{z}%
_{j}\left(  t\right)  \right\}
\end{equation}
which gives
\begin{equation}
\delta\mathbf{\rho}\left(  t\right)  =\left(
\begin{array}
[c]{cc}%
\frac{\partial\rho\left(  \mathbf{z}_{0},\mathbf{\bar{z}}_{0},t\right)
}{\partial\mathbf{\bar{z}}} & \frac{\partial\rho\left(  \mathbf{z}%
_{0},\mathbf{\bar{z}}_{0},t\right)  }{\partial\mathbf{z}}%
\end{array}
\right)  \left(
\begin{array}
[c]{c}%
\delta\mathbf{\bar{z}}\left(  t\right)  \\
\delta\mathbf{z}\left(  t\right)
\end{array}
\right)
\end{equation}


There is a linear relationship between the classical external field
$\mathcal{F}\left(  \mathbf{z},\mathbf{\bar{z}},t\right)  $ and the FORDM
$\mathbf{\rho}$ which is itself a linear function of the full density matrix
of the molecular system. Hence in the linear response regime where we only
consider first order changes in $\mathbf{\rho,}$ which are produced by
$\delta\mathbf{z}\left(  t\right)  ,$ it is consistent to approximate
$\mathcal{F}\left(  \mathbf{z+}\delta\mathbf{z},\mathbf{\bar{z}+}%
\delta\mathbf{\bar{z}},t\right)  $ about the point $\left(  \mathbf{z}%
_{0},\mathbf{\bar{z}}_{0}\right)  $ by
\begin{equation}
\mathcal{F}\left(  \mathbf{z}_{0}\mathbf{+}\delta\mathbf{z},\mathbf{\bar{z}%
}_{0}\mathbf{+}\delta\mathbf{\bar{z}},t\right)  \approx\mathcal{F}_{1}\left[
\mathbf{z}_{0}\right]  \left(  \delta\mathbf{z},\mathbf{\bar{z}}_{0},t\right)
=\mathcal{F}\left(  \mathbf{z}_{0},\mathbf{\bar{z}}_{0},t\right)
+D^{1}\mathcal{F}\left(  \mathbf{z}_{0},\mathbf{\bar{z}}_{0},t\right)  \left(
%
\begin{array}
[c]{c}%
\delta\mathbf{\bar{z}}\left(  t\right)  \\
\delta\mathbf{z}\left(  t\right)
\end{array}
\right)
\end{equation}
which produces the approximate field dependent classical Hamiltonian
\begin{equation}
\mathcal{E}_{2L}\left[  \mathbf{z}_{0}\right]  \left(  \delta\mathbf{z}%
_{0},\delta\mathbf{\bar{z}}_{0},t\right)  =\mathcal{E}_{0}\left[
\mathbf{z}_{0}\right]  \left(  \delta\mathbf{z}_{0},\delta\mathbf{\bar{z}}%
_{0}\right)  +\mathcal{F}_{1}\left[  \mathbf{z}_{0}\right]  \left(
\delta\mathbf{z}_{0},\delta\mathbf{\bar{z}}_{0},t\right)
\end{equation}



\end{document}